\section{Выводы}
В ходе выполнения практической работы было проведено имитационное моделирование сквозного цифрового тракта пространственно-временной обработки радиолокационных сигналов.

\begin{enumerate}
    \item \textbf{Работоспособность сквозного тракта.} Успешно реализована и отлажена цепочка преобразований, включающая аналого-цифровое преобразование, цифровой перенос частоты (ЦДМ), низкочастотную фильтрацию побочных составляющих (ЦФ) и согласованную фильтрацию (СФ).
    \item \textbf{Субдискретизация и перенос частоты.} Использование частоты дискретизации АЦП ($140$~МГц), которая меньше промежуточной частоты сигнала ($170$~МГц), на практике подтвердило работоспособность метода субдискретизации (bandpass sampling).
    \item \textbf{Вычислительная эффективность и тождественность методов.} Сравнительное моделирование подтвердило строгую математическую тождественность алгоритмов прямой свертки (FIR), спектрального метода (FFT/IFFT) и блочной сегментированной свертки (\texttt{fftfilt}). Метод \texttt{fftfilt} на основе алгоритма перекрытия со сложением является наиболее предпочтительным для обработки длинных потоков данных в реальном времени.
    \item \textbf{Влияние оконного взвешивания.} Применение весового окна Хэмминга к опорному сигналу позволило снизить уровень первого бокового лепестка с теоретических $-13.3$~дБ до $-42.8$~дБ. Данное подавление сопровождается закономерным расширением главного лепестка сжатого импульса по уровню $-3$~дБ примерно в $1.35 \dots 1.4$ раза.
\end{enumerate}